\section{边界上的交换：从军镇到关市}

\subsection{守边的人先看见粮车}

宋金边界在条约中可以用山河与州县标出，在地面上却由堡寨、渡口、烽火、田地和道路组成。守边士兵所等待的不只是敌军，也包括迟到的粮车、替换的同袍和来自后方的命令；边民所面对的不只是旗帜，更是军队借用牲畜、关卡检查货物和农时被警报打断。边防的稳定因而不能只从几年间没有大战判断。它需要一套持续运作的供给、侦察与地方协商。

绍兴、隆兴以后的和议给双方提供了交聘和停战的框架，却没有取消边境的摩擦。使节往来以国书、礼物和称谓安排两朝如何相见，边地交易则要处理盐、布帛、牲畜、药材与日常用品如何过关。朝廷的禁令能说明担忧什么，不能说明所有货物从此停止流动。沿边市场可能使居民获得稀缺物资，也可能成为军需、走私和情报交织的地点；同一条道路在和平时运货，在战时又会运兵。

\subsection{关市的税和风险}

关市不是一扇只为征税而开的门。官员需要辨别何物可过、何人可留、何时应禁；商人需要判断路途、关税、价格和战事消息；地方人则可能借熟悉的渡口与亲属网络避开或承受检查。我们没有金代每一关市的逐日账册，无法写出一张精确的贸易地图。正史中的朝贡与交聘条目也不应被误作市场统计。它们最有价值之处，是提示外交秩序、财政收入和地方生活在边界上并非互不相关。

边市交易同样不能被浪漫化为跨政权的自由市场。军队可能优先占用道路和货物，战时禁令会使小商人承担更大风险，地方豪强或官吏也可能借检查获利。即使一项交易有利于双方居民，也未必符合中枢的安全计算。正因如此，市场的存在并不反驳边界的暴力，而是说明居民在暴力的阴影中仍要寻找盐、粮、布和消息。

\subsection{淮河与长江的不同尺度}

淮河地带的渡口、湖泊和堤防使进攻与防御都依赖水陆配合。金军要南下，不能只依靠北方的骑兵机动；船只、纤夫、岸上仓储和水情都会限制推进。南宋的防守也不是天然得来的，它需要修堤、巡江、征粮和地方官把水网转化为可用的屏障。采石等战事因此应当被理解为长期后勤条件在特定地点的集中，而非一位将领或一种武器单独决定一切。

长江以北与以南的经济条件不同，战争却会把二者紧密相连。北方缺粮或道路受阻，可能增加对南方边界的压力；南宋的财政调度与江防，也会影响金廷判断进攻还是议和。诗歌中的北望和朝廷文书中的战守，都是这条边界的不同语言。它们可以说明人们如何感受等待与危险，不能替我们量出每座寨堡的兵额或每个家庭的忠诚。

\subsection{军镇里的社会}

军镇并非只有士兵。军人的家属、供应粮草者、铁匠、车夫、商贩、僧道和被临时征调的农户，都在其中留下或不留下痕迹。军籍把某些人归入可调用的队伍，实际生活却可能跨越军、民、商和亲属关系。退役、迁调、婚姻和市场交易会改变一户人在军镇中的位置。以军户或猛安谋克一个名称概括所有人，容易把这种流动消失。

边区材料保存尤其不均。城址能够显示防御空间，碑刻偶尔记录某次修建或某位将领，史书则记大战与调兵；它们都不够让我们重建完整的军镇人口。应当避免将一处关隘的遗迹推广为整个边防，也避免以史书沉默断言边民没有活动。可确认的是，边区社会承担了和平与战争的双重成本：既要适应检查、驻军与税课，也要在警报来临时承担更大的风险。

\subsection{外交中的称谓和现实}

兄弟、伯侄、臣属或和好等称谓，是外交文书用来组织权力关系的语言。它们的重要性不在于揭示一个永恒的上下等级，而在于规定使节可以怎样说话、礼物如何交付、边界争端以何种名义被提出。金与南宋都可在不同场合强调或淡化这些称谓，以应对各自的中枢政治与边防需要。将其中一个词直接翻译为现代主权关系，会抹去它在仪式和谈判中的策略性。

条约也不能自动证明地方服从。中央承认某条边界，仍要有士兵守路、官员征税、居民耕作和商人运输，才能使它成为日常可见的事实。某处争执、逃亡或私贸不必然推翻整项和议，却显示边界总是在执行中被重新定义。金朝长期维持南方关系的能力，正在这种看似琐碎的持续劳动中；后来压力累积时，首先断裂的也常是这些劳动所依赖的道路与信任。
